% !PRINTED 17
\weberpage{17}
%<<CONT
\[
 x+(-x)=0,
 \qquad x+0=x,
\]
\[
 x+(-y)=x-y,\quad \text{if }x>y,
 \quad =-(y-x),\quad \text{if }y>x.
\]
With this definition of addition, if $z_1,z_2,z_3$ are any three
numbers of the whole number domain, the laws
\[
 z_1+z_2=z_2+z_1,
 \qquad (z_1+z_2)+z_3=z_1+(z_2+z_3),
\]
hold; they are called the commutative and associative laws. The sum of
any number of summands is formed by successively combining any two
summands into a sum. Subtraction no longer needs separate
consideration if $z_1-z_2$ is defined as $z_1+(-z_2)$ and
$-(-z)=z$ is set.

The number line is represented intuitively by points: from a fixed
point of a straight line, denoted by $0$, the positive numbers are laid
off as segments on one side, say the right, and the negative numbers
on the other (left) side. The image of the sum of two segments
$z_1+z_2$ is obtained by laying off from the point $z_1$ the segment
of length $\pm z_2$ to the right or left, according as $z_2$ is
positive or negative.

Multiplication and division are introduced into the enlarged number
domain by the equations
\[
 x(-y)=(-x)y=-xy,
\]
\[
 (-x)(-y)=xy,
 \qquad 0x=0
\]
are defined. Division, as the inverse of multiplication, is always possible except
when the divisor is zero.

\begin{center}
\rule{0.18\linewidth}{0.4pt}
\end{center}

A further extension of the concept of number consists in introducing
\gesperrt{complex quantities}. We combine every two numbers of the
complete number line into pairs $(x,y)$, and regard two such pairs
$(x,y)$ and $(a,b)$ as equal only when $x=a$, $y=b$. These pairs of
numbers form a manifold whose elements cannot be ordered, but with
which addition, subtraction, multiplication, and division are to be
performed according to the following rules. Let
\[
 (x,y)+(a,b)=(x+a,y+b),
\]
\[
 (x,y)(a,b)=(xa-yb,xb+ya),
\]
%CONT>>
